\documentclass[11pt,letterpaper]{article}
\usepackage[T1]{fontenc}
\usepackage[utf8]{inputenc}
\usepackage{amsmath,amsthm,mathtools}
\usepackage{newtxtext,newtxmath}
\usepackage[margin=1in,headheight=15pt]{geometry}
\usepackage{microtype,booktabs,tabularx,longtable,array}
\usepackage{enumitem,algorithm,algpseudocode}
\usepackage{xcolor,fancyhdr,titlesec,needspace}
\usepackage[round,authoryear]{natbib}
\usepackage{xurl,hyperref,listings}
\definecolor{ink}{HTML}{302447}
\definecolor{muted}{HTML}{625970}
\hypersetup{colorlinks=true,linkcolor=ink,citecolor=ink,urlcolor=ink,
 pdftitle={Earn the Next Experiment: Transformational Intelligence through an Improvement-for-Capacity Pilot},
 pdfauthor={Vincent Boucher},pdfsubject={A proposed AGI Alpha Agent / SUCCESSOR Omega pilot with executed reference studies},
 pdfkeywords={transformational intelligence, improvement-for-capacity, SUCCESSOR Omega, resource acquisition, causal evaluation}}
\urlstyle{same}
\pagestyle{fancy}\fancyhf{}
\fancyhead[L]{\small\textsc{Earn the Next Experiment}}
\fancyhead[R]{\small Boucher}
\fancyfoot[C]{\thepage}
\renewcommand{\headrulewidth}{0.35pt}
\setlength{\parindent}{1.1em}\setlength{\parskip}{0.10em}
\setlength{\emergencystretch}{2em}
\setlength{\tabcolsep}{5pt}\renewcommand{\arraystretch}{1.12}
\titleformat{\section}{\large\bfseries\color{ink}}{\thesection}{0.65em}{}
\titleformat{\subsection}{\normalsize\bfseries}{\thesubsection}{0.65em}{}
\titlespacing*{\section}{0pt}{1.20em}{0.5em}
\titlespacing*{\subsection}{0pt}{0.75em}{0.25em}
\setlist{nosep,leftmargin=1.5em,topsep=.2em}
\renewcommand{\topfraction}{.95}\renewcommand{\bottomfraction}{.85}
\renewcommand{\textfraction}{.05}\renewcommand{\floatpagefraction}{.85}
\newcolumntype{Y}{>{\raggedright\arraybackslash}X}
\newcolumntype{P}[1]{>{\raggedright\arraybackslash}p{#1}}
\newcommand{\I}{\mathcal I}\newcommand{\A}{\mathcal A}\newcommand{\E}{\mathbb E}
\newcommand{\Prob}{\mathbb P}\newcommand{\J}{\mathcal J}\newcommand{\1}{\mathbf 1}
\DeclareMathOperator{\Feas}{Feas}
\newcommand{\thesisbox}[1]{\begin{center}\setlength{\fboxsep}{9pt}\fcolorbox{ink}{white}{\begin{minipage}{.92\linewidth}#1\end{minipage}}\end{center}}
\newtheorem{definition}{Definition}[section]
\newtheorem{proposition}{Proposition}[section]
\newtheorem{theorem}{Theorem}[section]
\theoremstyle{remark}\newtheorem{remark}{Remark}[section]
\lstset{basicstyle=\small\ttfamily,breaklines=true,columns=fullflexible,keepspaces=true,
 showstringspaces=false,frame=single,rulecolor=\color{muted},aboveskip=.65em,belowskip=.65em}
\begin{document}
\thispagestyle{plain}
\begin{center}
{\fontsize{29}{33}\selectfont\bfseries\color{ink} Earn the Next Experiment\par}
\vspace{.6em}
{\Large\bfseries Transformational Intelligence through an\par}
\vspace{.12em}
{\Large\bfseries Improvement-for-Capacity Pilot\par}
\vspace{1em}
{\large Vincent Boucher\par}
{MONTREAL.AI\par}
{Montr\'eal, Qu\'ebec, Canada\par}
\vspace{.6em}
{October 2026\quad$\cdot$\quad Research manuscript\par}
\end{center}
\begin{abstract}
An AI can use a fixed computing budget more efficiently while remaining dependent on someone else to finance its next experiment. We study a different formulation: improve a willing provider's computational workload, establish the benefit independently, receive a bounded allocation of usable service capacity as compensation, and use that delivered capacity on a second mission. The proposed \emph{improvement-for-capacity} pilot makes resource acquisition an outcome of useful intelligence rather than an assumed input or a token-price forecast. Its central distinction is between a faster method, released resources, a contractual entitlement, actual delivery, and a useful successor result. We define time- and resource-specific capacity, give a conditional non-displacement certificate and a two-sided participation interval, and specify a complete pilot inside AGI Alpha Agent / SUCCESSOR $\Omega$. Two executed reference studies expose the decisive failure modes. In a finite allocation model, a validated twelve-unit gain earns three service credits, but equal aggregate capacity can either complete a three-slot descendant task or fail because the available intervals are fragmented. Across all 8,092 supplied delivery profiles with the same aggregate release, 666 cannot serve that task. A separate repeated-linear-system experiment confirms that retained factorization beats repeated reconstruction on the tested batched inputs, while an ordinary batched library baseline is better still. Neither laboratory is an external provider pilot or autonomous algorithm-discovery result. The contribution is a practical, falsifiable path from a discovered transformation to an additional resource under usable control, with full costs, provider obligations, capacity expiry, and the next mission kept visible. The next experiment must actually run; a speedup, credit, or valuation is not its substitute.
\end{abstract}
\noindent\textbf{Keywords:} transformational intelligence; improvement-for-capacity; resource acquisition; capability construction; causal evaluation; SUCCESSOR $\Omega$.
\begin{center}\small
\textbf{Find your Alpha. Build your SUCCESSOR $\Omega$.}\par
\href{https://github.com/MontrealAI/AGI-Alpha-Agent-v0}{github.com/MontrealAI/AGI-Alpha-Agent-v0}
\end{center}
\clearpage

\section{The next resource need not be purchased in advance}
The conventional question is: \emph{How can we obtain more compute to make our AI more useful?} An alternative is: \emph{Whose existing system can our AI make more productive, and would that provider compensate a verified improvement with capacity we can actually use?} The second question changes the production and compensation arrangement, not just the speed of a computation.

The proposed pilot has one provider, one public-input primary workload, one independent evaluator, and one separately specified descendant mission. Alpha investigates the provider's restriction. A qualified improvement preserves the provider's obligations. The provider retains most of the verified benefit and fulfills a bounded service allocation under agreed terms. SUCCESSOR $\Omega$ uses that allocation to complete the descendant mission. Failure at any link remains a visible result.

\thesisbox{\textbf{Transformational intelligence recognizes when the original problem formulation is unnecessarily restrictive---and discovers a fundamentally different game to play.}\par\smallskip Here the changed game is not merely spending a fixed budget better. It is earning an additional productive resource by making another system more productive.}

This is not computation from nothing. Hardware, energy, labor, contractual performance, and initial research resources still have owners and costs. The provider must be better off relative to a credible alternative, and the network must receive something usable rather than a nominal credit. The scheme can fail even after a technically impressive discovery.

\subsection{What this paper adds to the project brief}
The accompanying project brief proposes a provider whose validated monthly output rises from 100,000 to 110,000 units, with 7,500 of the additional units retained by the provider and 2,500 allocated to the network. Those numbers are an illustrative offer, not an observed business result \citep{brief}. We preserve its core pilot: a useful improvement, an independently checked exchange, and a second mission completed with the delivered capacity.

The paper adds the missing conversion discipline. Units of the first workload do not automatically become units of the second. An improvement can reduce one bottleneck while leaving the required memory, time window, software environment, or permission unavailable. A provider's future-service promise is also not stored physical resource: yesterday's unused processor time cannot be delivered tomorrow. A new service must be supplied at the time it is owed.

Our preceding manuscripts developed formulation change and productive-capacity formation \citep{boucherOriginal,boucherProduction}. The present increment narrows the empirical target to one observable exchange. It does not claim that its first cycle is already self-funding, that a token appreciates, or that repeated expansion follows indefinitely. It asks a testable question: \emph{did the useful transformation cause the system to receive and successfully use a resource it did not otherwise have under the compared arrangement?}

\subsection{Contributions and claim boundary}
The contribution is an operational research design with four parts: a capacity-conversion model that separates five evidential claims; a time- and resource-specific delivery certificate with explicit counterexamples; a prospective provider/evaluator/operator protocol; and executable reference artifacts, including a numerical-workload control and an exhaustive finite allocation study.

Performance contracting, abstraction, scheduling, and AI-assisted optimization have substantial precedents. Energy-savings performance contracts already link compensation to improvements and require measurement and verification \citep{doeContract,doeMV}. AlphaEvolve already connects model-generated programs to measured computational improvements \citep{alphaevolve}. We do not claim that payment in kind or factorization reuse is new. The research target is their disciplined composition into an AI capability-acquisition experiment, with a demonstrated descendant outcome rather than a narrated feedback loop.

\section{Five different claims between discovery and a successor}
\label{sec:claims}
A useful pilot cannot collapse all progress into ``capacity gained.'' Table~\ref{tab:claims} separates the observable objects. The first two can hold without a commercial agreement; the agreement can exist without fulfillment; fulfillment can occur without a useful descendant result.

\begin{table}[htbp]\centering\small
\caption{A chain of evidence, not interchangeable success labels.}\label{tab:claims}
\begin{tabularx}{\linewidth}{@{}P{.17\linewidth}YY@{}}
\toprule
Claim & Required observation & What it does not establish\\\midrule
Technical improvement & Required outputs and service conditions preserved, with a measured resource or performance gain & Transferable capacity or payment\\
Net resource release & Gain survives implementation, verification, monitoring, and displaced-resource costs & That the provider owes or can deliver it\\
Earned entitlement & Authorized parties accept a bounded service obligation under explicit terms & Delivered processor time or cash\\
Capacity delivered & The specified environment supplies the authorized resource window to the successor & Useful output, full cost recovery, or token demand\\
Descendant benefit & The second mission produces an accepted artifact using that delivery & Better discovery across campaigns or indefinite growth\\\bottomrule
\end{tabularx}
\end{table}

\subsection{The invariant is the provider's purpose, not its old workflow}
The primary contract specifies work that must be completed, output quality, permissible delay, data rights, resource limits, and prohibited side effects. The AI may challenge an inherited decomposition, repeated preprocessing, a representation, a scheduling dependency, or a measurement procedure. It may not remove required outputs or unilaterally relax quality to create a favorable resource ledger.

For example, independent linear-system requests can share a coefficient matrix. A workflow may refactor that matrix for each request because requests were treated as isolated units. A maintained representation may remove that restriction. But waiting to form a large batch is not legitimate when it violates a required response deadline. Material changes to the service contract require an amendment and a new comparison basis, not retroactive improvement of an earlier score.

\subsection{The creative target includes the exchange}
A technically useful transformation need not yield the kind of capacity the descendant needs. Therefore Alpha should search jointly over the method, an informative experiment, a proposed compensation arrangement, and a feasible descendant mission. Let a candidate be
\begin{equation}
T=(F_T,\pi_T,e_T,k_T,B_T),
\end{equation}
where $F_T$ is the formulation, $\pi_T$ the implementation, $e_T$ the test, $k_T$ the proposed service terms, and $B_T$ the descendant plan. This is a proposal space, not a delegation of contract authority. The controller approves the relevant terms before execution or external commitment.

Joint search need not change the intended scientific mission to whatever is cheapest. If the desired descendant is fixed, the exchange must serve it. If the user has authorized a portfolio of useful descendants, the system may select among them transparently. Replacing a demanding required mission with an easy demonstration is not a successful conversion.

\section{Capacity is a profile, not a scalar balance}
\label{sec:capacity}
\subsection{Typed resource accounting}
Define a delivery environment by hardware class $h$, software image $v$, location $\ell$, time slots $t\in\{1,\ldots,H\}$, and resources $r\in\mathcal R$. Resources may include processor occupancy, accelerator memory, host memory, bandwidth, storage, and energy. Their units remain separate. Let $C_{rt}$ be the physical service limit, $P_{rt}$ the preserved provider workload, $O_{rt}$ incremental overhead, and $Z_{rt}$ previously committed service. A new allocation $x_{rt}$ must satisfy
\begin{equation}\label{eq:physical}
P_{rt}+O_{rt}+Z_{rt}+x_{rt}\le C_{rt}\quad\text{for every }r,t.
\end{equation}
Eligibility also requires the specified environment, rights, delivery interval, and service quality. Multi-resource allocation is a standard systems problem; scalar CPU totals alone do not resolve it \citep{ghodsi,borg}.

For a primary metering resource, let $L_{rt}$ denote an admitted lower envelope of incremental release relative to the agreed counterfactual, after the relevant overhead. Where a pilot promises to allocate only improvement-origin resources, it additionally requires
\begin{equation}\label{eq:incremental}
x_{rt}+Z^{\mathrm{improvement}}_{rt}\le L_{rt}.
\end{equation}
An instantaneous lower envelope is a stronger object than a lower bound on average savings. When the available evidence supports only an average, the provider must separately back each delivery reservation with actual headroom; it cannot convert an expectation into a guaranteed per-slot release.

Auxiliary resources can be supplied from existing provider capacity, but their provenance, cost, and authorization remain explicit. CPU improvement does not create GPU memory. The reference study uses existing authorized memory headroom alongside newly available CPU service; it does not count that memory as created by the improvement.

\subsection{Earned service credits are claims on a future service}
After a primary performance tranche is measured and accepted, a simple illustrative contract awards
\begin{equation}\label{eq:credit}
q=\min\{q_{\max},\lfloor\alpha G\rfloor\},\qquad 0\le\alpha\le1,
\end{equation}
where $G\ge0$ is qualified net improvement in a stated metering unit. The conversion from this unit to a future service unit must be a contractual term with an appropriate calibration. It is not a physical identity. The credit names a provider, recipient, class of service, expiry, and limit; it is not a freely circulating financial asset in the proposed pilot.

The agreement can be negotiated before discovery, but an entitlement is marked earned only after its stated conditions hold. A contingent future reservation can be made earlier at a provider's explicitly accepted risk; it remains a contingent commitment rather than realized delivery. Unused, expired, revoked, disputed, or partially fulfilled credits must remain distinguishable from resources consumed by an accepted mission.

For each unit class, the operational ledger preserves
\begin{equation}\label{eq:balance}
q_{\mathrm{issued}}=q_{\mathrm{available}}+q_{\mathrm{reserved}}+q_{\mathrm{consumed}}+q_{\mathrm{expired/cancelled}}.
\end{equation}
The categories are disjoint. Cancellation does not erase an actual consumption. A completed request cannot be replayed as a new entitlement. A provider-authoritative ledger must enforce this transactionally across all clients; a stateless receipt checker cannot discover allocations absent from its input.

\subsection{A descendant-relative ordering}
Let $\Feas_{\J}(C)$ be the subset of a declared task family $\J$ that can be completed under profile $C$, including hardware, timing, resource, and permission constraints. Write $C_1\succeq_{\J}C_0$ when
\begin{equation}\label{eq:order}
\Feas_{\J}(C_0)\subseteq\Feas_{\J}(C_1).
\end{equation}
This definition makes the intended successor outcome part of the capacity comparison. It does not supply a tractable general scheduler or a universal price for resources. Two profiles with equal aggregate CPU time can be incomparable under different task families.

\begin{proposition}[Equal totals need not buy the next experiment]\label{prop:fragment}
There exist equal-total released-resource profiles and equal earned credits such that a nonpreemptible descendant task is feasible in one profile and infeasible in the other.
\end{proposition}
\begin{proof}
Over eight equal time slots, take available CPU profiles
\[
g=(2,2,2,2,1,1,1,1),\qquad g'=(3,0,3,0,3,0,3,0).
\]
Both total twelve. A three-credit task requiring one CPU unit in each of three consecutive slots fits in $g$. Every positive interval in $g'$ has length one, so it cannot fit there. Assume adequate memory and the same admitted hardware and window. The impossibility comes from temporal structure, not total quantity.
\end{proof}

\thesisbox{\textbf{Do not optimize only how much capacity an improvement appears to release. Optimize whether the released and contracted service can run the next valuable experiment.}}

\subsection{A conditional non-displacement certificate}
A certificate supplies the frozen primary contract and outcome checks, the measurement basis, a bounded authorized entitlement, the current calendar, all relevant existing commitments, and a descendant schedule. A separate checker validates Equations~\eqref{eq:physical}--\eqref{eq:balance}, environment compatibility, expiry, task completeness, and the descendant's result.

\begin{theorem}[Finite-scope delivery and non-displacement]\label{thm:certificate}
Assume the input records accurately describe the provider's obligations, resource limits, existing commitments, and authorized service terms. If the proposed schedule satisfies every per-resource, per-slot constraint, consumes no more than the earned entitlement, completes the descendant's required operations within its valid window, and both primary and descendant output checks pass, then the modeled schedule completes the descendant without violating the represented provider obligations or reusing a represented unit of committed capacity.
\end{theorem}
\begin{proof}
The resource inequalities exclude overload in each time slot and resource dimension. The schedule and window checks establish the required descendant operations. The ledger identity and unique reservation identities exclude duplicate consumption of represented credit. The output checks establish the stipulated outcomes. Applying these facts to every slot gives the claimed finite-horizon result.
\end{proof}

The theorem is a consistency certificate under supplied premises, not a proof of honest metering, complete demand, correct hardware, legal enforceability, or future provider solvency. Missing obligations and forged measurements can invalidate the premises. Independent evidence custody and operational observation remain necessary. The result is not new scheduling theory; its purpose is to specify exactly what the claimed exchange must demonstrate.

\section{Two participants must prefer the exchange}
\label{sec:economics}
\subsection{A transparent participation interval}
Consider a simplified homogeneous-unit model. Let $\Delta>0$ be qualified released service units over a declared horizon, $\alpha$ the contracted network share, $v_P,v_N>0$ the respective net values per retained or usable unit, and $d\in(0,1]$ the fraction of awarded units that become useful network service under the assumed delivery model. Let $K_P,K_N\ge0$ include all participant costs not already included in these per-unit values. Define
\begin{align}
U_P(\alpha)&=(1-\alpha)\Delta v_P-K_P,\\
U_N(\alpha)&=d\alpha\Delta v_N-K_N.
\end{align}
No party's token holdings are included as operating returns. The provider's values include opportunity cost, including possible sale of the released capacity elsewhere. The network's values reflect real useful service, not the provider's list price by assumption.

\begin{proposition}[Nonnegative participation interval]\label{prop:participation}
Under this linear model, both parties have nonnegative net value exactly when
\begin{equation}\label{eq:interval}
\max\left\{0,\frac{K_N}{d\Delta v_N}\right\}
\le\alpha\le
\min\left\{1,1-\frac{K_P}{\Delta v_P}\right\}.
\end{equation}
A mutually nonnegative exchange exists if and only if the interval is nonempty.
\end{proposition}
\begin{proof}
Rearrange $U_N\ge0$ and $U_P\ge0$ and intersect with $[0,1]$.
\end{proof}

This is an elementary necessary-and-sufficient calculation for a supplied linear valuation model, not an incentive-compatible mechanism or an empirical estimate. Strict preference requires strict inequalities. Outside a risk-neutral expectation model, uncertain delivery needs a more appropriate risk criterion. Values, costs, bargaining power, and alternatives are not known merely because the equation is available.

For illustration, use $\Delta=12$, $v_P=2$, $v_N=3$, and $K_P=K_N=6$. With $d=1$, the feasible share interval is $[1/6,3/4]$. At $d=1/2$, it narrows to $[1/3,3/4]$; at $d=1/4$, to $[2/3,3/4]$. At $d=1/8$, no share works. A nominally attractive quarter-share is not robust to these delivery assumptions. These are hypothetical common-value units, not prices for any real provider.

\subsection{A seed is allowed; calling it recovered requires evidence}
The first experiment needs initial resources. Include model calls, failed candidates, engineering, benchmark execution, independent review, contract work, installation, and monitoring. Those resources may be financed by the founder, a sponsor, or a provider. They must not be omitted to manufacture a self-funding claim.

Three outcomes should be reported separately. \emph{Mechanism success} means useful work earned and redeemed a bounded service allocation. \emph{Resource recovery} means the relevant earned resources cover the declared seed resource vector under justified conversion rules. \emph{Commercial viability} means the complete arrangement creates acceptable value after all costs and obligations. One can hold without the others.

The exact study intentionally assigns five seed CPU units and delivers only three modeled CPU units. It therefore demonstrates no recovery of that seed, even in its successful delivery case. Cash and human effort are unmeasured. Initial losses can be legitimate research investment; they are not evidence of autonomous financial sustainability.

\subsection{Do not induce the wrong improvement}
An agreement based on gross CPU time can encourage displaced memory costs, batching that violates latency, inflated baselines, or optimizations a provider would have installed anyway. The counterfactual, horizon, exclusions, and attribution rule must precede paid construction. A meaningful benchmark contains a strong existing alternative, not only the historical workflow.

Compensation should attach to the accepted service benefit. Review compensation is fixed or otherwise independent of a positive verdict, subject to separately agreed terms. The evaluator must be able to reject a gain, a credit calculation, or delivery without losing earned compensation for performing the review. Performance contracting's measurement literature provides useful precedents for allocating such risks \citep{doeMV,doeWitness}; this manuscript does not claim their legal templates automatically apply to computing providers.

\section{The first natural pilot: a public numerical service}
\label{sec:pilot}
\subsection{Selection criteria before a provider is named}
The initial workload should be material to a willing provider, reproducible on public or synthetic nonpersonal inputs, and independently checkable. It should have a plausible removable restriction and a descendant workload compatible with the likely freed resource. It must not need customer credentials, private repositories, production writes, or a speculative hardware purchase to establish the first result.

We propose a repeated linear-system service as an initial task class, not a claim that this is the best available commercial opportunity. The provider must first show that the service exists, that its cost matters, and that a strong readily available implementation has not already captured the proposed benefit. If those conditions fail, choose a different public numerical or software workload rather than preserve the demonstration for appearances.

\subsection{Primary task A and its possible reformulation}
For a supplied symmetric positive-definite matrix $A$ and multiple right-hand sides $b_i$, return solutions to
\begin{equation}
Ax_i=b_i,\qquad i=1,\ldots,m,
\end{equation}
under declared numerical and latency conditions. An isolated-request formulation repeats matrix factorization for every $b_i$. A purpose-directed formulation recognizes that the matrix, not the request boundary, determines the reusable preprocessing. It can retain a factorization or use a batched solve where the contract permits it. LAPACK explicitly separates factorization from solves using the resulting factors, and SciPy exposes that distinction \citep{lapackFactor,lapackSolve,scipyFactor,scipySolve}.

This is an established numerical technique. Its recovery by an AI on unfamiliar source could be useful without being a scientific invention. The local experiment supplies the methods; it does not demonstrate that Alpha independently discovered them. In a field pilot, the investigator records the model's starting information, source access, proposals, rejected candidates, tool calls, and actual modifications. The system is not handed the intended winning reformulation and then credited with discovering it.

A matrix identity must bind actual content, shape, numeric type, and relevant environment. Reusing a factor after the matrix changes is not a valid speedup. Preprocessing and identity checks are charged. The admissible design must also account for factor memory, lifetime, contention, and recovery. An isolated cache keyed only by a convenient name is not sufficient.

\subsection{Acceptance checks and a strong comparator}
The evaluator checks every required output, input preservation, finite values, and a declared backward-error criterion, such as
\begin{equation}\label{eq:residual}
\frac{\|Ax_i-b_i\|_{\infty}}
{\|A\|_{\infty}\|x_i\|_{\infty}+\|b_i\|_{\infty}}
\le\varepsilon.
\end{equation}
A residual is not automatically a forward-error guarantee on an ill-conditioned problem. Generated reference cases can include known solutions for additional checks; real inputs require condition-sensitive requirements. Symmetry and positive-definiteness must be part of task admission, not inferred from a successful call on one selected matrix triangle.

Compare the incumbent with a tuned library baseline, a strong unrestricted agent with the same tools and budget, and the proposed method. A known library implementation that wins should be adopted or retained when it serves the provider. No mechanism requires an LLM-branded candidate to win. Novelty, autonomous discovery, measured improvement, and commercial relevance are separate claims.

\subsection{Choose task B before claiming useful capacity}
The descendant is selected with an explicit purpose, data, checker, hardware class, software image, memory needs, input/output limits, timing, and an execution budget. Suitable first missions include a fresh public numerical investigation or a bounded resource-allocation computation with an independent oracle. It should produce a useful artifact, not merely print a success message or rerun the exact primary benchmark.

The strongest pilot selects B independently of the eventual speedup. If B is a new experiment in a capability-development campaign, a valid negative scientific result may still be useful. More computation does not guarantee a new discovery. Receipt of capacity, completion of the experiment, and improvement in the successor must each retain their own status.

\section{Measurement must earn the right to support a service obligation}
\label{sec:measurement}
\subsection{Choose one primary accounting basis}
Two natural comparisons answer different questions. A \emph{fixed-work} experiment measures resources needed to deliver the same workload. A \emph{fixed-resource-window} experiment measures validated throughput under identical resource and service limits. Either can be appropriate. Paying once for saved time and again for the extra tasks that occupy the same saved time double-counts one benefit.

The first pilot should use fixed work with a declared service horizon, because the required output set can be compared directly. It should report processor occupancy, elapsed time, memory, data movement, and energy when measured. Wall-clock savings on a partly idle service are not automatically transferable core-seconds. Likewise, a lower kernel time does not prove a lower end-to-end cost.

\subsection{Development and confirmation are different epochs}
Development may be adaptive. Confirmation uses a frozen candidate, checker, software and hardware specification, selected endpoints, and fresh input batches or service blocks. Record all failed candidates and their cost. If a new candidate is selected after seeing confirmatory results, it begins another accounted attempt. Repeated use of a holdout can compromise ordinary inference \citep{dwork}.

For a resettable isolated numerical service, use paired baseline/candidate runs with randomized order, equivalent initialization, and no shared candidate cache leaking into the baseline. Preserve both cold-start and intended steady-state comparisons. A real shared service may require an appropriately designed switchback experiment; carryover, traffic cycles, and interference must be treated explicitly, not ignored because the same machine is used \citep{switchback}.

\subsection{A conservative qualification bound is not deliverable inventory}
For illustration, let $D_i\in[-M,M]$ be independent paired net savings observations for a frozen method, with common expectation $\mu$. Include the stipulated overhead in $D_i$. Hoeffding's inequality gives
\begin{equation}\label{eq:lcb}
L_n=\overline D_n-M\sqrt{\frac{2\log(1/\beta)}{n}},\qquad
\Prob(\mu\ge L_n)\ge1-\beta.
\end{equation}
This follows by applying the bounded-difference inequality with range width $2M$ \citep{hoeffding}. A practical gain requires $L_n$ above a prespecified materiality threshold as well as acceptable quality. Estimated, unbounded, dependent, or nonstationary observations do not satisfy this statement merely because they are entered into the formula.

Allocate error probability across repeated confirmatory attempts, for example $\beta_j=\beta/[j(j+1)]$ on fresh appropriately sampled data. If continuous monitoring is desired, use a valid sequential method rather than repeatedly inspecting a fixed-sample interval \citep{howard}. Neither construction repairs a biased measurement instrument.

Most importantly, Equation~\eqref{eq:lcb} bounds an average under its sampling assumptions. It does not reserve a future CPU, promise a memory envelope, or ensure delivery before expiry. The provider must separately accept the obligation, make an exclusive reservation, and provide observed execution. Statistical evidence qualifies a method; the scheduler and operating contract qualify the resource transfer.

\subsection{Provider and evaluator responsibilities}
The provider supplies truthful permitted baseline information, controls the execution environment, and honors accepted service obligations. The evaluator controls confirmation inputs and the acceptance calculation and checks provider-level output preservation. Alpha constructs candidate methods within its admitted tools. The owner authorizes resource use and contracts. These may be separate roles without being independent organizations; the field report must identify actual control and conflicts rather than count addresses or agent names.

Metering evidence should be available to the evaluator without publishing personal information or private operational data. The pilot's public-input boundary concerns the computational material. Necessary business contact and contract administration, if any, remain with the responsible parties rather than being scraped into the discovery model.

\section{A complete proposed operating protocol}
\label{sec:protocol}
\subsection{An agreement small enough to execute and specific enough to check}
The project brief's 25\% network share is an example, not a fixed optimum. Negotiate a share and cap only after the provider's alternative, resource calibration, downside, and delivery terms are understood. A research term sheet should specify Table~\ref{tab:terms}. It is a technical specification for counsel and counterparties, not an executable legal agreement.

\begin{table}[htbp]\centering\small
\caption{Minimum proposed terms for one provider and one descendant.}\label{tab:terms}
\begin{tabularx}{\linewidth}{@{}P{.24\linewidth}Y@{}}\toprule
Term & Required content\\\midrule
Primary outcome & Exact workload, quality, deadlines, rights, prohibited effects, and rollback criteria\\
Counterfactual & Frozen incumbent, tuned alternatives, adaptation allowance, metering unit, attribution horizon\\
Research budget & Seed compute, cash, evaluator fees, operator time, failed attempts, maximum exposure\\
Compensation & Earned-gain definition, share, cap, conversion to future service, maturity and dispute terms\\
Delivery envelope & Provider, hardware, image, region, time window, CPU/accelerator and memory limits, auxiliary costs\\
Reservation rules & Priority, exclusivity, expiry, concurrency, cancellation, outage response, unused-credit treatment\\
Descendant purpose & Named mission B, permitted tools, acceptance evidence, recovery and terminal classification\\
Rights and publication & Artifact ownership or license, open-source obligations, allowed evidence publication\\
Review & Independent scope, compensation including valid negative findings, escalation and conflicts\\
Economic separation & No implicit company equity, token-holder ownership, customer-principal reuse, or mandatory native payment\\\bottomrule
\end{tabularx}
\end{table}

\subsection{The two-epoch sequence}
Epoch A closes only after the primary improvement has been measured and reviewed. Its accepted evidence establishes whether a bounded service entitlement has been earned. Epoch B has later resource windows and must be backed by the provider's actual service inventory. The resource consumed in epoch B is not the same physical time that went unused in epoch A. Ongoing improved operation may make the later delivery affordable, but that is an operating obligation to observe, not a conservation identity across time.

A provider can reserve an upcoming slot before every possible future uncertainty is resolved. The record must identify the remaining risk. A completed pilot records what was actually delivered; it never marks the full face value of unused credit as acquired useful capacity.

\begin{algorithm}[htbp]\small
\caption{One improvement-for-capacity pilot}
\begin{algorithmic}[1]
\Require Primary purpose A, intended descendant B, willing provider, evaluator, explicit budgets
\State Admit public inputs, rights, provider terms, and initial research resources
\State Reproduce incumbent and strongest relevant available alternative
\State Diagnose a removable restriction; propose formulations and discriminating tests
\State Construct and test admitted candidates; preserve failed attempts and costs
\State Freeze selected method, source, checker, calibration, and confirmation plan
\State Independently assess quality, net gain, provider obligations, and materiality
\If{no qualified gain or no provider acceptance}
  \State Retain incumbent; close as rejected or inconclusive; do not issue earned credit
\Else
  \State Close epoch A; record capped earned service entitlement
  \State Validate an exclusive later resource reservation for B
  \If{reservation or execution fails}
    \State Preserve the unfulfilled/partial obligation; apply agreed recovery or expiry
  \Else
    \State Execute B through the actual provider environment; meter actual use
    \State Independently check B; reconcile consumption and remaining credit
    \State Retain artifact and evidence under one SUCCESSOR project identity
  \EndIf
\EndIf
\State Report mechanism success, seed recovery, and economic viability separately
\end{algorithmic}
\end{algorithm}

\subsection{The exact stop conditions}
Stop or renegotiate when baseline reproduction fails, a stronger readily available method removes the proposed advantage, quality degrades, costs cannot be bounded, permission is missing, useful delivery cannot fit before expiry, or the provider cannot honor reservations without displacing represented obligations. A lost sponsor budget is not a reason to relabel expected gains as receipts. A negative result can be accepted research work under a separately funded agreement without generating performance-based credit.

Interruptions preserve actual consumed resources and outstanding liabilities. Retries need a new or explicitly remaining budget; they must not restore previously consumed credit. A changed checker or purpose requires separately reviewed comparison terms. A changed provider environment requires requalification appropriate to its effect.

\section{Reference study I: execute the candidate workload, keep the stronger control}
\label{sec:numeric}
\subsection{What was actually run}
The companion runs three supplied methods on synthetic positive-definite systems: refactor for each right-hand side; retain one factorization and solve columns separately; and use one factorization with a batched library solve. The retained-factor method also computes a content fingerprint within its timed call. The matrix and right-hand sides are generated as
\begin{equation}
A=R^\top R+nI,\qquad B=AX_\star,
\end{equation}
with fixed seeded pseudorandom $R,X_\star$. This produces controlled reference systems, not naturally collected provider workloads. Each output is checked by Equation~\eqref{eq:residual} with $\varepsilon=10^{-12}$ and a reference forward-error bound of $10^{-10}$.

The recorded environment was Python 3.13.5, NumPy 2.3.5, SciPy 1.17.0, Linux x86-64. BLAS thread environment variables were set to one before import. No dedicated-host, thermal, or CPU-isolation claim is made. Each case has one warmup per method and 21 cyclic-order timed repetitions. Input generation and result checking are excluded equally; each timed invocation rebuilds the needed factor. Raw nanosecond samples are retained. These timings are descriptive observations on one environment, not confidence bounds on a provider population.

\begin{table}[htbp]\centering\small
\caption{Measured median elapsed milliseconds for the complete supplied numerical calls. No result is converted into an earned provider credit.}\label{tab:numeric}
\begin{tabular}{@{}rrrrr@{}}\toprule
Matrix size & Right-hand sides & Refactor each & Retain factor & Batched library\\\midrule
192 & 1 & 0.188625 & 0.431742 & 0.183346\\
192 & 24 & 4.948795 & 1.383107 & 0.285521\\
384 & 24 & 23.352769 & 4.709395 & 1.245649\\\bottomrule
\end{tabular}
\end{table}

\subsection{The result that matters is not the largest ratio}
On the 192-by-192, 24-right-hand-side case, retained factorization is about 3.58 times faster than repeated reconstruction. Yet the ordinary batched library implementation is about 4.84 times faster than the retained-column method. At one right-hand side, the extra retained-method fingerprinting is overhead rather than an investment recovered by reuse. All recorded outputs passed the specified checks; the largest recorded reference forward error was below $10^{-14}$.

A pilot that advertised only its win over repeated reconstruction would miss the better available method. The correct conclusion is not that a special discovery architecture beats established numerical software. It is that the proposed workload and verification procedure can expose a removable request-boundary restriction, that reuse has a horizon-dependent cost, and that strong alternatives should determine selection.

A real provider may already use the batched implementation, may require online low-latency responses that preclude batching, or may be limited by another resource. In those cases this supplied candidate may create no incremental commercial value. The provider workload must be assessed before a performance agreement is offered. This experiment made no frontier-model call, contacted no capacity provider, and delivered no external computing allocation.

\section{Reference study II: an earned amount can fail to become useful capacity}
\label{sec:exact}
\subsection{A closed illustrative primary epoch}
The exact laboratory deliberately separates its virtual resource units from the milliseconds in Section~\ref{sec:numeric}. Its primary evidence record supplies a baseline of 32 reference CPU units, a candidate of 18, and two additional units of incremental overhead. Thus
\begin{equation}
G=32-18-2=12.
\end{equation}
Quality and confirmation are supplied true predicates in the positive case; they are not observations of a provider. A quarter share and a three-unit cap earn $q=3$ under the assumed one-to-one future-service conversion. The primary epoch closes at time 8; the entitlement is issued then and can be used only in later slots 9--16, with expiry at 17.

The provider's later calendar has four CPU and four memory units per slot. Its required CPU occupancy is $4-g_t$ and required memory occupancy is three. The incremental eligible CPU release is $g_t$. One memory unit per slot is already available and separately authorized; it is not charged as newly created. The descendant requires one CPU and one memory unit in three consecutive slots. Its entire modeled primary CPU credit must fit within the three-unit award.

\subsection{Aligned and fragmented delivery}
The aligned profile is $g=(2,2,2,2,1,1,1,1)$. The planner reserves slots 9, 10, and 11, and the separate checker validates all represented resource, timing, credit, and outcome predicates. The fragmented profile $g'=(3,0,3,0,3,0,3,0)$ has the same twelve-unit aggregate and three-credit entitlement, yet supplies no three-slot contiguous run. It therefore records zero delivered units and no descendant result.

\begin{table}[htbp]\centering\small
\caption{Exact supplied cases. ``Delivered'' refers only to the virtual schedule; no real provider exchange occurred.}\label{tab:cases}
\begin{tabularx}{\linewidth}{@{}Yrrr@{}}\toprule
Case & Raw net gain & Earned credit & Modeled delivery\\\midrule
Aligned resource window & 12 & 3 & 3\\
Equal-total fragmented window & 12 & 3 & 0\\
Memory unavailable & 12 & 3 & 0\\
Expired entitlement & 12 & 3 & 0\\
CPU service versus required GPU & 12 & 3 & 0\\
Primary quality regression & 12\textsuperscript{a} & 0 & 0\\
Unconfirmed improvement & 12\textsuperscript{a} & 0 & 0\\
Overhead erases gross gain & 0 & 0 & 0\\
Provider consent absent & 12\textsuperscript{b} & 0 & 0\\
Pre-existing idle time only & 0 & 0 & 0\\
Already committed release & 12 & 3 & 0\\
Fragmented, explicitly preemptible task & 12 & 3 & 3\\\bottomrule
\end{tabularx}
\par\smallskip\raggedright\footnotesize
\textsuperscript{a}The raw arithmetic difference remains 12, but quality or confirmation denies qualification. The JSON retains both the raw number and denial flags.\quad
\textsuperscript{b}Technical evidence alone does not establish an earned obligation.
\end{table}

The preemptible control allows the same three resource slots at times 9, 11, and 13. It is a separately declared task contract. It does not qualify the original nonpreemptible mission. A useful future discovery might construct a checkpointable method that preserves the original purpose, but its checkpoint overhead, memory, and outcome equivalence would need their own evidence. A flag change alone is not that discovery.

\subsection{Complete finite census and separate checking}
Enumerate all eight-slot profiles $g\in\{0,1,2,3\}^{8}$ with $\sum_tg_t=12$. There are 8,092 profiles. With the same credit and descendant requirement, 7,426 admit a schedule and 666 do not. The planner searches consecutive starts. A separately written oracle searches all three-slot subsets and checks adjacency and capacity. They agree for every profile, and every generated schedule passes the separate certificate checker.

This is exhaustive enumeration of a specified finite space, not 8,092 independent AI discovery trials. The infeasible fraction is a property of that enumeration measure; it is not an estimate of real provider failure rates. Its purpose is to demonstrate why an aggregate-capacity headline can conceal a completely unusable reward.

The companion also includes a single-process credit-book reference. It reserves before consumption, rejects overspending and conflicting operation identities, records actual partial use, and makes identical terminal updates idempotent. A concurrent test admits only one of eight three-unit reservation requests against a three-unit balance. This is not a distributed provider ledger, a persistent crash-safe service, or blockchain qualification.

\subsection{A real local computation under a virtual allocation}
On the positive schedule, the driver executes a separately declared eight-item finite allocation task using ordinary dynamic programming. A brute-force checker examines all 256 subsets, establishes feasibility, and confirms the optimum value of 27 under the supplied budget of 17. The resulting selection is retained as a descendant artifact.

The arithmetic execution is real; the service allocation and three-slot cost are modeled. The code does not sleep for three provider time slots, constrain its process to an external allocation, or prove that the allocated resource physically paid for that computation. This distinction is explicit in the artifact. The complete external pilot remains unperformed.

The laboratory additionally records five seed CPU units. Delivering three modeled units cannot establish seed recovery. That negative accounting result prevents a single successful exchange-shaped computation from being advertised as a self-sustaining growth cycle.

\section{Implementation inside AGI Alpha Agent / SUCCESSOR \texorpdfstring{$\Omega$}{Omega}}
\label{sec:implementation}
\subsection{Extend the project rather than invent a parallel platform}
The named implementation target is
\begin{center}
\url{https://github.com/MontrealAI/AGI-Alpha-Agent-v0}\par\smallskip
\textbf{Find your Alpha. Build your SUCCESSOR $\Omega$.}
\end{center}
The main ref inspected for this manuscript resolves to \texttt{b44ba296e00500c91c1617fc942b5f0050ceabd6}. Its flagship guide describes public-input mission preparation, competing strategies, native Computer Use, retained returns, and owner review. It explicitly distinguishes proposals from executed patches or measured savings, and states that arbitrary code execution and payments are outside that browser adapter \citep{alphaGuide}. These are snapshot facts, not a claim that every project branch has the same scope.

The improvement-for-capacity path is a proposed additive increment. It does not exist merely because the guide can display the plan. It requires a qualified numerical execution adapter, provider-reservation integration, metering, independent review, and actual recipient-side execution. A browser can assist with permitted investigation; it does not have to imitate clicks when an authorized computational interface better serves the purpose.

\subsection{One identity across two missions and one exchange}
Retain a parent project identifier and distinct identities for primary A, improvement T, confirmation epoch, service agreement, earned entitlement, reservation, descendant B, and terminal evidence. Reuse existing journal and project-history infrastructure where available. Source hashes bind inputs and outputs; they do not authenticate economic independence or prove the truth of provider telemetry.

The proposed records are intentionally typed. An \texttt{ImprovementEvidence} identifies the comparison and quality verdict. A \texttt{ServiceEntitlement} identifies what a provider owes after acceptance. A \texttt{CapacityReservation} identifies a currently admitted resource window. A \texttt{ConsumptionReceipt} identifies actual use. A \texttt{DescendantOutcome} identifies the resulting artifact and review. These may be linked but cannot replace one another.

\begin{table}[htbp]\centering\small
\caption{Proposed additive engineering boundaries; file names are suggestions, not existing-feature claims.}\label{tab:modules}
\begin{tabularx}{\linewidth}{@{}P{.38\linewidth}Y@{}}\toprule
Proposed component & Responsibility\\\midrule
\texttt{capacity\_pilot.py} & Mission pair, lifecycle transitions, explicit permissions and stop conditions\\
\texttt{capacity\_measurement.py} & Counterfactual, outcome evidence, overhead, calibration and qualified gain\\
\texttt{capacity\_agreement.py} & Nonbinding preparation and separately accepted provider terms\\
\texttt{capacity\_reservation.py} & Provider API adapter, exclusive windows, expiry, idempotent reservation and reconciliation\\
\texttt{capacity\_delivery.py} & Actual installed-recipient execution, resource receipts, partial failure and descendant binding\\
\texttt{capacity\_pilot\_view} & English/French preparation, evidence, actual costs, remaining obligations and portable history\\\bottomrule
\end{tabularx}
\end{table}

Native schemas must be versioned. An old zero-allowance request must not acquire spending or external-service authority through reinterpretation. Imported histories are evidence to inspect, not current provider approval. The service credential remains with the authorized operator; no private key, API secret, or usable authorization is embedded in a public artifact.

\subsection{The smallest release-worthy increment}
Build one complete path on the current coordinated development line after its actual source and interfaces are selected. The scope is one provider adapter and one computational task class, not arbitrary cloud control. Required engineering evidence includes valid reservation and refusal cases, exact environment binding, interruption after partial use, replay prevention, expiry, output recovery, and a complete installed-package journey on the supported platform.

Preserve ongoing security corrections, existing review gates, historical releases, original papers and URLs, deployed contracts, ENS settings, and the USDC/\$AGIALPHA route distinction. The present manuscript performs no repository write, merge, release, deployment, service purchase, or token action. A source release would itself establish neither provider commissioning nor a completed economic pilot.

\subsection{A user interface that exposes the decisive state}
The user should see a small set of plain-language states: \emph{proposed improvement}, \emph{benefit independently checked}, \emph{provider owes a specified service}, \emph{capacity reserved}, \emph{capacity actually used}, and \emph{second mission reviewed}. Show unknown costs and pending obligations rather than a single celebratory progress indicator.

The final card should answer four questions: What restriction changed? What did the provider retain? What resource was actually delivered? What useful result did the successor produce with it? The English/French journey must preserve the same underlying meanings and exact source identities. A visible percentage gain is secondary to those answers.

\section{The \$AGIALPHA connection is a separate empirical step}
\label{sec:token}
The first provider exchange need not involve a blockchain payment. Compensation in service capacity can be a bilateral operating arrangement with an ordinary off-chain provider account. Do not create a second transferable token merely to name a nontransferable service obligation.

The project brief proposes using some earned capacity for paid services through an admitted $\$\mathrm{AGIALPHA}$ route \citep{brief}. That creates a testable subsequent hypothesis: an independent customer buys a useful service, the service runs on capacity actually delivered under the provider agreement, and the native settlement is independently reconciled with the artifact and full costs. It is not necessary for the first technical exchange to prove that hypothesis.

Record at least three ledgers separately: the provider's primary service benefit; the network's earned and consumed capacity, seed costs, and obligations; and any actual native-token receipts, fees, or cash conversion. A token transfer alone does not demonstrate external demand. A capacity credit is not cash, and unused face value is not realized profit. A successful company or provider network does not automatically grant token holders rights to its profits, equipment, or intellectual property.

No token appreciation, price floor, passive yield, company equity, or ownership of provider infrastructure is promised by this design. Any new financial entitlement or token offering needs its own legal and economic analysis; Canadian guidance expressly addresses offerings described as utility tokens \citep{csa}. Existing commissioning and settlement rules remain distinct, and no live financial transaction is performed by the reference artifact.

The long-horizon connection is practical rather than promotional: useful discoveries may earn resources; useful resource deliveries may make additional experiments affordable; some later services may attract genuine native commercial demand. Each transition is separately measured. The first exchange is evidence for one link, not an estimate of the entire future economy.

\section{How the pilot can fail---and still advance the research}
\label{sec:evaluation}
\subsection{Six falsifiable hypotheses}
\paragraph{H1: useful formulation discovery.}
On an independently selected public-input workload, a configured Alpha agent discovers and realizes a change that improves the agreed outcome or full resource cost against strong same-permission alternatives. Supplying the winning method demonstrates realization, not autonomous discovery.

\paragraph{H2: capacity conversion.}
The admitted gain supports a bounded entitlement that the provider fulfills in a compatible environment without displacing represented primary obligations. A speedup with no accepted terms, an unredeemed credit, or unusable intervals does not support this claim.

\paragraph{H3: descendant usefulness.}
The delivered service executes a separately specified mission whose artifact passes independent review. Replaying A, generating a trivial success marker, or changing B's requirements after seeing the available resources does not establish the intended benefit.

\paragraph{H4: bilateral economic acceptability.}
Both parties prefer the exchange to their credible alternatives after complete costs and opportunity costs. A provider's courtesy credit, sponsor subsidy, waived evaluator fee, or unresolved support liability must be disclosed rather than counted as recurring earned surplus.

\paragraph{H5: valid reuse and recovery.}
The retained improvement remains useful on new compatible work and is declined or requalified when its assumptions change. Stale matrix identities, provider-image drift, missed updates, expiry, and failed partial consumption are explicit tests, not excluded inconvenience cases.

\paragraph{H6: improved acquisition strategy.}
A learned policy becomes better at choosing improvements, experiments, exchange terms, and descendants over separate complete campaigns under matched starting resources. One successful pilot, a larger inherited capacity balance, or many tasks within one campaign does not establish improved discovery skill.

\subsection{A factorial that distinguishes technique from the exchange}
Cross technical improvement absent/present with capacity-compensation agreement absent/present. All cells have the same externally authorized opportunity to negotiate and use the tools; cells are analytical treatments, not unequal permissions hidden from a baseline. Without improvement there may be no earned compensation. Without accepted exchange terms there may be no new resource under network control. The combined cell tests the proposed mechanism.

Also compare buying an equivalent service with cash, payment in cash for the same improvement, and a strong provider-internal optimizer. Capacity compensation is not universally preferable. The provider may value cash differently, the network may need a different resource class, or contracting may cost more than the allocation is worth. A real negotiated comparison is stronger evidence than a synthetic financial model.

\subsection{Report the complete stopping point}
The field report should name the strongest completed claim from Table~\ref{tab:claims}, then report quality, resources, terms, delivery, descendant output, and complete economics. Avoid a single aggregate ``pilot passed'' label when one transition is unresolved. Record provider and reviewer independence accurately and keep any security or legal qualification separate from numerical correctness.

A single natural mission supports an instance result. Fresh missions test repeatability. General improvement in resource acquisition requires separately initialized campaigns with held-out terminal evaluations. Metareasoning and automated agent design provide relevant comparisons for allocating search effort and learning procedures \citep{hay,adas}; a new name for the loop does not establish superiority.

\section{Relationship to prior work and limits of the present evidence}
\subsection{What is inherited and what is new here}
Our formulation-change framework is continuous with the earlier project papers \citep{boucherOriginal,boucherProduction}. Program abstraction and reusable languages have technical precedent in DreamCoder \citep{dreamcoder}; AI-assisted program improvement has direct precedent in AlphaEvolve \citep{alphaevolve}. This manuscript adds a narrowly specified conversion from an independently accepted improvement to actual recipient-side capacity and a separately checked descendant result.

Payment from verified improvement is not new. Energy performance contracting offers a close conceptual precedent, although its regulatory, asset, and financing conditions differ from a computing-service pilot \citep{doeContract,doeMV}. Advance market commitments likewise demonstrate that demand can be organized before a desired capability is available \citep{amc}. The pilot does not inherit the empirical success of either precedent.

Resource profiles and fragmentation belong to established scheduling and multi-resource systems work \citep{ghodsi,borg}. Our finite propositions do not replace those methods. They locate the failure that a scalar ``compute earned'' claim overlooks. The separate receipt checker borrows the general engineering separation of production and verification; it remains conditional on trusted inputs, implementation, and custody.

\subsection{Limitations}
No independent capacity provider was recruited for this manuscript. No external contract was signed, no capacity was earned or redeemed from a third party, and no native-token settlement occurred. The local numerical methods and exact planning procedures were investigator-supplied. The workload and schedule families are synthetic, fully specified, and selected to expose mechanisms; they are not demonstrations of general transformational intelligence.

The numerical benchmark measures elapsed time in one local environment, not hardware capacity, energy savings, customer value, or an enforceable service claim. The exact resource model uses artificial units and supplied true/false outcome predicates. Its descendant arithmetic is actually executed but not physically funded or metered by the virtual allocation. The separate checker shares authorship and premises with the constructor and is not an independent institution or a general security boundary.

The contract analysis assumes simplified values and delivery fractions. It omits the richer bargaining, risk, tax, legal, market, and operational analysis needed for a real agreement. Efficient hardware use can increase subsequent resource consumption; no environmental benefit follows without its own measurement. Successful resource acquisition may still fail to improve a scientific result, repay the seed budget, or support a durable business.

\section{Conclusion: make the next experiment an earned outcome}
The most consequential change is not a larger claim about how much intelligence might eventually be worth. It is a different immediate experiment: use intelligence to remove a provider's costly restriction, let the provider retain a demonstrable benefit, and receive a service that can actually support the next valuable mission.

The reference studies make two requirements visible. Strong existing methods must be allowed to beat a favored AI candidate. And an identical aggregate gain can have completely different usefulness depending on when, where, and with which complementary resources it is delivered. The successful unit is not the percentage speedup or the balance of credits. It is the preserved primary service plus the actually enabled descendant outcome.

For AGI Alpha Agent, \emph{Find your Alpha} means finding a restriction whose removal matters to a real beneficiary. \emph{Build your SUCCESSOR $\Omega$} means retaining the useful capability, the evidence, and a legitimately earned resource that supports another mission. One completed provider exchange would be a modest but decisive empirical step. It would not prove unlimited compounding; it would establish a link that ambitious compounding stories often assume away.

\thesisbox{\textbf{Do not merely ask for more resources so that intelligence can become useful. Test whether useful intelligence can earn the resources for what it should do next.}}

\section*{Research status and reproducibility}
This manuscript provides a proposed natural pilot, analytic arguments, an executed local numerical benchmark, an exact allocation census, a separate consistency checker, and a local credit-book reference. All 64 tests pass in the recorded environment. The associated project brief supplies the strategic premise; the models, protocols, calculations, and implementations here extend it. External sources support technical and institutional context. The original papers and project repositories are unchanged.

\clearpage\appendix
\section{A ready-to-fill pilot charter}
The companion includes the following nonbinding preparation structure as a JSON template. All absent provider, authorization, and evidence identities remain null. A template is not a signed agreement or spending authority.
\begin{lstlisting}
{
  "schema": "ifc-pilot-charter-proposal/v1",
  "status": "UNCOMMISSIONED",
  "project": "AGI Alpha Agent / SUCCESSOR Omega",
  "provider": null,
  "independent_evaluator": null,
  "primary_workload": {
    "type": "public numerical service",
    "source_and_input_manifest": null,
    "required_outputs_and_quality": null,
    "strong_baselines": null,
    "resource_boundary_and_service_level": null
  },
  "compensation": {
    "kind": "bounded nontransferable service allocation",
    "share": null,
    "cap_and_unit": null,
    "calibration": null,
    "delivery_window": null,
    "provider_acceptance": null
  },
  "descendant": {
    "purpose": null,
    "environment_and_resources": null,
    "independent_checker": null
  },
  "seed_costs": null,
  "approvals": null,
  "token_settlement": "none by default"
}
\end{lstlisting}

\subsection{Provider invitation}
The first outreach can be specific without promising a breakthrough:
\begin{quote}
We propose one bounded experiment on a public-input computational workload. We will reproduce your current service and a strong available alternative, investigate a potentially removable restriction, and commission independent review of any candidate improvement. Before work starts, we agree on quality, resource accounting, exposure, and a capped service-capacity allocation that would be earned only if the specified benefit is accepted. You retain the negotiated majority of the benefit. We then test whether the delivered allocation completes a separate useful mission. No result, payment, or capacity gain is presumed in advance.
\end{quote}
The responsible parties must supply actual budgets, legal terms, provider identity, and review arrangements. No outreach or transaction is sent by the companion.

\section{Reproduction, files, and evidence interpretation}
\subsection{Commands}
The exact allocation study and its checker use only Python's standard library. The numerical workload additionally uses NumPy and SciPy. Recorded versions are Python 3.13.5, NumPy 2.3.5, and SciPy 1.17.0.
\begin{lstlisting}
python experiments/capacity_pilot.py --output reproduced_exact
python experiments/verify_capacity.py reproduced_exact/allocation_cases.json
python -m unittest discover -s tests -v
# Optional: preserves the supplied timing observations.
python experiments/numerical_workload.py --output reproduced_benchmark
\end{lstlisting}
The numerical command requires a fresh output directory and preserves the supplied observations. Timings are not expected to reproduce byte-for-byte. The exact allocation JSON files are deterministic and can be compared by hash after rerunning. The top-level \texttt{reproduce.py} runs the exact study in a temporary directory, checks byte agreement, validates its certificates, and executes the tests.

Compile the standalone manuscript with \texttt{latexmk -pdf Earn\_the\_Next\_Experiment.tex}, or run \texttt{pdflatex} twice. Its bibliography and tables are embedded; there are no external figure or font files. The package includes the original brief as a preserved source note, the proposed charter, implementation handoff, result files, and claim register.

\begin{table}[htbp]\centering\small
\begin{tabularx}{\linewidth}{@{}P{.45\linewidth}Y@{}}\toprule
File & Scope\\\midrule
\texttt{experiments/capacity\_pilot.py} & Supplied cases, bounded schedule search, exact census, participation arithmetic, and child task\\
\texttt{experiments/verify\_capacity.py} & Independently written finite-record checks and brute-force child oracle\\
\texttt{experiments/credit\_book.py} & Single-process reservation, consumption, idempotence, and conservation reference\\
\texttt{experiments/numerical\_workload.py} & Three supplied numerical methods and local timing/output checks\\
\texttt{results/allocation\_cases.json} & Twelve complete exact-case certificates\\
\texttt{results/allocation\_summary.json} & Case summary, census and illustrative participation intervals\\
\texttt{results/numerical\_benchmark.json} & Recorded environment, raw timings, factors and numerical checks\\
\texttt{tests/} & 64 accounting, scheduling, output, mutation and concurrency tests\\
\texttt{CLAIMS.md} & What each result supports and what remains unperformed\\\bottomrule
\end{tabularx}
\end{table}

The consistency checker accepts bounded local inert data. It does not validate arbitrary code, act as a sandbox, authenticate a real provider, or protect a shared cluster. Its nonfinite/duplicate-key and identity checks are ordinary integrity tests of the supplied record shape. Distributed ledger durability, actual metering, and external review remain separate implementation work.

\subsection{Precise separation of the two studies}
No measured millisecond is converted to a virtual CPU-slot in the exact allocator. The twelve-unit primary gain is a supplied scenario input, not derived from a measured numerical call. Conversely, the real numerical outputs establish neither an earned quota nor a transfer of service ownership. The pair of studies tests different components of the pilot design; it does not secretly substitute for their external integration.

The live referee-controlled campaign, credible provider workload, agreed compensation, normal installed-provider adapter, actual resource reservation, observed consumption, and independent beneficiary acceptance are the missing empirical link. Their absence is visible in the charter and scope fields rather than represented by synthetic identities that look live.

\section{Claim--evidence register}
\small
\begin{longtable}{@{}P{.31\linewidth}P{.33\linewidth}P{.30\linewidth}@{}}\toprule
Claim & Evidence & Limit\\\midrule\endfirsthead
\toprule Claim & Evidence & Limit\\\midrule\endhead
Useful intelligence can be investigated as a source of new resource access & Proposed two-epoch pilot and explicit exchange terms & Hypothesis until an external provider performs\\
Equal aggregate capacity can differ in descendant usefulness & Proposition~\ref{prop:fragment}; two exact profiles & Specified nonpreemptible task and resource class\\
A checked schedule preserves represented obligations & Theorem~\ref{thm:certificate}; separate checker & Conditional on truthful complete records and execution\\
A mutually acceptable share may not exist & Proposition~\ref{prop:participation}; exact fractions & Supplied linear values, costs and delivery assumptions\\
Retained factorization beats repeated reconstruction in tested batched cases & Local raw timings and output checks & Supplied algorithms; batched library control is faster\\
7,426 of 8,092 profiles admit the specified task & Complete exact enumeration and separate oracle & Not a population rate or independent campaigns\\
A descendant allocation result is actually calculated & Dynamic programming and all-256-subsets checker & Computation is real; allocation and its charge are virtual\\
Credit cannot be double spent within the tested local book & Reservation/terminal/concurrency tests & Single process, no persistence or external authority\\
The positive exact case recovers the seed budget & Not supported: five seed units exceed three delivered units & Cash and human effort unmeasured\\
AGI Alpha Agent is the implementation target & Pinned flagship guide and proposed adapter map & No integration, merge, release, or activation here\\
Native token demand or valuation rises & Not established & Separate commercial test, no price or ownership promise\\
The strategy improves its own resource-acquisition skill & H6 and matched full-campaign design & No completed learned-strategy superiority study\\\bottomrule
\end{longtable}
\normalsize

\clearpage
\begin{thebibliography}{99}
\footnotesize\setlength{\itemsep}{0pt}\setlength{\parsep}{0pt}\setlength{\parskip}{0pt}
\bibitem[Bojinov et~al.(2020)]{switchback}
Iavor Bojinov, David Simchi-Levi, and Jinglong Zhao. \emph{Design and Analysis of Switchback Experiments}. arXiv:2009.00148, 2020. \url{https://arxiv.org/abs/2009.00148}.
\bibitem[Boucher(2026a)]{boucherOriginal}
Vincent Boucher. \emph{Beyond the Original Problem: Transformational Intelligence for Expanding What Is Possible}. MONTREAL.AI, October 2026. \url{https://github.com/MontrealAI/AGI-Alpha-Agent-v0/blob/main/docs/Beyond_the_Original_Problem.pdf}.
\bibitem[Boucher(2026b)]{boucherProduction}
Vincent Boucher. \emph{Beyond the Production Frontier: Transformational Intelligence from Verified Discovery to Stellar-Scale Civilization}. MONTREAL.AI, October 2026. \url{https://github.com/MontrealAI/AGI-Alpha-Agent-v0/blob/main/docs/Beyond_the_Production_Frontier.pdf}.
\bibitem[Canadian Securities Administrators(2018)]{csa}
Canadian Securities Administrators. \emph{Canadian securities regulators provide additional guidance on securities law implications for offerings of tokens}. June 11, 2018. Accessed October 2026. \url{https://www.securities-administrators.ca/news/canadian-securities-regulators-provide-additional-guidance-on-securities-law-implications-for-offerings-of-tokens/}.
\bibitem[Dwork et~al.(2014)]{dwork}
Cynthia Dwork, Vitaly Feldman, Moritz Hardt, Toniann Pitassi, Omer Reingold, and Aaron Roth. \emph{Preserving Statistical Validity in Adaptive Data Analysis}. arXiv:1411.2664, 2014. \url{https://arxiv.org/abs/1411.2664}.
\bibitem[Ellis et~al.(2020)]{dreamcoder}
Kevin Ellis et~al. \emph{DreamCoder: Growing generalizable, interpretable knowledge with wake-sleep Bayesian program learning}. arXiv:2006.08381, 2020. \url{https://arxiv.org/abs/2006.08381}.
\bibitem[Ghodsi et~al.(2011)]{ghodsi}
Ali Ghodsi, Matei Zaharia, Benjamin Hindman, Andy Konwinski, Scott Shenker, and Ion Stoica. Dominant Resource Fairness: Fair Allocation of Multiple Resource Types. In \emph{NSDI}, 2011. \url{https://www.usenix.org/conference/nsdi11/dominant-resource-fairness-fair-allocation-multiple-resource-types}.
\bibitem[Hay et~al.(2012)]{hay}
Nicholas Hay, Stuart Russell, David Tolpin, and Solomon Eyal Shimony. \emph{Selecting Computations: Theory and Applications}. arXiv:1207.5879, 2012. \url{https://arxiv.org/abs/1207.5879}.
\bibitem[Hoeffding(1963)]{hoeffding}
Wassily Hoeffding. Probability inequalities for sums of bounded random variables. \emph{Journal of the American Statistical Association}, 58(301):13--30, 1963. \url{https://doi.org/10.1080/01621459.1963.10500830}.
\bibitem[Howard et~al.(2021)]{howard}
Steven R. Howard, Aaditya Ramdas, Jon McAuliffe, and Jasjeet Sekhon. Time-uniform, nonparametric, nonasymptotic confidence sequences. \emph{The Annals of Statistics}, 49(2):1055--1080, 2021. \url{https://arxiv.org/abs/1810.08240}.
\bibitem[Hu et~al.(2024)]{adas}
Shengran Hu, Cong Lu, and Jeff Clune. \emph{Automated Design of Agentic Systems}. arXiv:2408.08435, 2024. \url{https://arxiv.org/abs/2408.08435}.
\bibitem[Kremer et~al.(2020)]{amc}
Michael Kremer, Jonathan Levin, and Christopher M. Snyder. Advance Market Commitments: Insights from Theory and Experience. \emph{AEA Papers and Proceedings}, 110:269--273, 2020. \url{https://doi.org/10.1257/pandp.20201017}.
\bibitem[LAPACK(2026a)]{lapackFactor}
LAPACK. \emph{POTRF: triangular factor}. Official reference documentation, accessed October 2026. \url{https://www.netlib.org/lapack/explore-html/d2/d09/group__potrf.html}.
\bibitem[LAPACK(2026b)]{lapackSolve}
LAPACK. \emph{POTRS: triangular solve using factor}. Official reference documentation, accessed October 2026. \url{https://www.netlib.org/lapack/explore-html/d3/dc8/group__potrs.html}.
\bibitem[MONTREAL.AI(2026a)]{brief}
MONTREAL.AI. \emph{Improvement-for-capacity strategy and pilot brief}. User-supplied working text, October 2026; attachment \texttt{Pasted text(20261020-002306).txt}, preserved in the companion source notes. Strategic proposal, not a field-result report.
\bibitem[MONTREAL.AI(2026b)]{alphaGuide}
MONTREAL.AI. \emph{Useful Alpha: a capability worth keeping}. AGI Alpha Agent, inspected commit \texttt{b44ba296e005}\ldots, October 2026; exact source below. \url{https://github.com/MontrealAI/AGI-Alpha-Agent-v0/blob/b44ba296e00500c91c1617fc942b5f0050ceabd6/docs/agent/SUCCESSOR_ALPHA.md}.
\bibitem[Novikov et~al.(2025)]{alphaevolve}
Alexander Novikov et~al. \emph{AlphaEvolve: A coding agent for scientific and algorithmic discovery}. arXiv:2506.13131, 2025. \url{https://arxiv.org/abs/2506.13131}.
\bibitem[SciPy(2026a)]{scipyFactor}
SciPy. \emph{scipy.linalg.cho\_factor}. Official API documentation. Recorded experiment uses SciPy 1.17.0. \url{https://docs.scipy.org/doc/scipy-1.17.0/reference/generated/scipy.linalg.cho_factor.html}.
\bibitem[SciPy(2026b)]{scipySolve}
SciPy. \emph{scipy.linalg.cho\_solve}. Version 1.17.0 API documentation. \url{https://docs.scipy.org/doc/scipy-1.17.0/reference/generated/scipy.linalg.cho_solve.html}.
\bibitem[U.S. Department of Energy(2026a)]{doeContract}
U.S. Department of Energy, Federal Energy Management Program. \emph{Energy Savings Performance Contracts for Federal Agencies}. Official program description, accessed October 2026. \url{https://www.energy.gov/cmei/femp/energy-savings-performance-contracts-federal-agencies}.
\bibitem[U.S. Department of Energy(2026b)]{doeMV}
U.S. Department of Energy, Federal Energy Management Program. \emph{Measurement and Verification for Federal Energy Savings Performance Contracts}. Accessed October 2026. \url{https://www.energy.gov/cmei/femp/measurement-and-verification-federal-energy-savings-performance-contracts}.
\bibitem[U.S. Department of Energy(2026c)]{doeWitness}
U.S. Department of Energy, Federal Energy Management Program. \emph{Guide to Government Witnessing and Review of Measurement and Verification Activities}. Accessed October 2026. \url{https://www.energy.gov/cmei/femp/articles/guide-government-witnessing-and-review-measurement-and-verification-activities}.
\bibitem[Verma et~al.(2015)]{borg}
Abhishek Verma, Luis Pedrosa, Madhukar Korupolu, David Oppenheimer, Eric Tune, and John Wilkes. Large-scale cluster management at Google with Borg. In \emph{EuroSys}, 2015. \url{https://research.google/pubs/large-scale-cluster-management-at-google-with-borg/}.
\end{thebibliography}
\end{document}
